\section{Experiments}
\label{sec:experiments}

We report four experiments, in the order of the questions in
Section~\ref{sec:intro}. Section~\ref{sec:compare} compares all the models,
Section~\ref{sec:crop} tests the tactile crop, Section~\ref{sec:sensors} asks
which sensors ACT needs \new{and which inputs each model uses}, and
Section~\ref{sec:wm} scores the world action models' predictions of the
cameras.

\subsection{Comparing the models}
\label{sec:compare}

Table~\ref{tab:actions} gives every model's action error on the training and
test demonstrations, and Figure~\ref{fig:horizon} shows how the error grows
with how far ahead the model plans.

\begin{table}[t]
  \centering
  \caption{Action error (RMSE, degrees) on the 58 training and the 7 test
  demonstrations. The left block scores each model over its own plan; the
  right block scores every model over the first ten steps
  (\SI{0.33}{\second}), the only block that is fair across rows. \emph{Gap} is
  test error divided by training error. \new{The lowest error in each column is
  in bold.}}
  \label{tab:actions}
  \small
  \begin{tabular}{lrrrrrr}
\toprule
& & \multicolumn{3}{c}{over its own plan} & \multicolumn{2}{c}{over the first 10 steps} \\
\cmidrule(lr){3-5}\cmidrule(lr){6-7}
model & plan & train & test & gap & test & gap \\
\midrule
ACT & 100 & 2.67 & 14.09 & 5.3 & 7.88 & 3.4 \\
Diffusion & 32 & 4.74 & 9.18 & 1.9 & 5.03 & 2.8 \\
pi0.5 & 50 & 17.05 & 18.35 & 1.1 & 8.96 & 1.1 \\
pi0.5, three passes & -- & -- & -- & -- & -- & -- \\
Flow matching & 50 & 7.57 & 14.44 & 1.9 & 4.51 & 1.4 \\
DreamZero & 24 & \textbf{0.94} & 7.02 & 7.4 & 4.17 & 5.3 \\
FastWAM & 10 & 1.53 & \textbf{3.94} & 2.6 & \textbf{3.94} & 2.6 \\
\bottomrule
\end{tabular}

\end{table}

\begin{figure}[t]
  \centering
  \includegraphics[width=0.72\textwidth]{horizon_all.png}
  \caption{Action error at each step of the plan, against how far ahead that
  step is, for every model on one linear scale. Left: test demonstrations.
  Right: training demonstrations. The shaded band is the first ten steps, which
  every model plans.}
  \label{fig:horizon}
\end{figure}

\textbf{Both world action models plan the next ten steps better than every
other model on unseen demonstrations.} Over those steps their test errors are
3.9 (FastWAM) and 4.2 (DreamZero). The two are too close to rank. Their lead is
several times larger than the change we see when we rerun them with different
random seeds (Appendix Table~\ref{tab:floor}). It does not prove that
predicting the future is what helps: FastWAM starts from a pretrained video
model, and DreamZero \new{is given the past \SI{1.6}{\second} of arm commands}.

\rev{\textbf{Flow matching is the most accurate behaviour-cloning policy}, over
the first ten steps and further ahead: 4.3 against Diffusion's 5.0 over ten
steps, 7.3 against 7.6 over 24 and 9.0 against 9.2 over 32. Flow matching and
pi0.5 share their training objective, and flow matching's error over the first
ten steps is half pi0.5's (4.3 against 9.0). On our data, pi0.5's pretraining
with our narrow LoRA does worse than training a small model from scratch. Our
first flow-matching run saw only a random \(84\times84\) patch of each camera
image, about 2\,\% of it (Section~\ref{sec:models}), and 98.7\,\% of what moved
its plan was proprioception. Retrained on whole,
resized images, it improves on the test demonstrations from 14.4 to 11.5 over
its plan, but only from 4.5 to 4.3 over the first ten steps. Over a third of a
second the arms' own positions predict the next commands almost entirely;
seeing the cameras matters further ahead.}

\textbf{Diffusion's error starts low but grows fastest.} On the test
demonstrations it rises from 3.9 at the first step to 13.9 by step 32, a factor
of 3.6; ACT starts at 7.3 but grows by only 1.8, so the two cross at about one
second (Figure~\ref{fig:horizon}). On the training demonstrations ACT's error is
almost flat over all 100 steps (2.4 to 3.3): it replays a familiar
demonstration almost exactly.

\textbf{Every model is much worse on unseen demonstrations.} ACT's error is
about five times higher on the test demonstrations than on the training ones,
and Diffusion's about twice as high. \new{Because the last seven demonstrations
are also the latest, a slow change over the afternoon could make them harder.
When the seven are instead picked at random, each sits between training
demonstrations in time, so drift should make them easier; the gap grows instead
(8.9 against 5.3 for ACT, 2.6 against 1.9 for Diffusion; Appendix
Figure~\ref{fig:split}). Drift does not explain the gap.}

\textbf{pi0.5's small gap does not mean it generalises well.} Its error on
the demonstrations it trained on is already higher than Diffusion's on
demonstrations it never saw: after about one pass over the data, adapting only
the small LoRA addition, it never fitted the training set. \new{Trained for
about three passes (\num{70000} steps), it barely moved: training error fell
from 17.1 to 16.4 and test error from 18.4 to 18.0, less than rescoring with
another seed moves it (Appendix Table~\ref{tab:floor}). The narrow LoRA, not
too little training, is the likelier limit.}

\rev{\subsection{Do the test demonstrations look different?}
\label{sec:recordings}

Figure~\ref{fig:recordings} shows how far apart every pair of the 65
demonstrations looks to each model's own vision encoder
(Section~\ref{sec:measures}), and Table~\ref{tab:recordings} asks the question
that matters for the gap above: is a test demonstration further from its
nearest training demonstration than a training demonstration is from its
nearest neighbour?

\begin{figure}[t]
  \centering
  \includegraphics[width=0.85\textwidth]{recordings.png}
  \caption{How different every pair of demonstrations looks to each model's
  vision encoder: the mean feature distance after matching their frames in
  order, scaled so that 1 is the most different pair. Demonstrations are in
  recording order; the lines mark the seven test demonstrations (58--65).}
  \label{fig:recordings}
\end{figure}

\begin{table}[t]
  \centering
  \caption{Distance from each demonstration to its nearest training
  demonstration (itself excluded), averaged over the 58 training and the 7 test
  demonstrations, on the scale of Figure~\ref{fig:recordings}.}
  \label{tab:recordings}
  \small
  \begin{tabular}{lrrr}
\toprule
model & training & test & test / training \\
\midrule
ACT & 0.26 & 0.28 & 1.07 \\
Diffusion & 0.36 & 0.32 & 0.90 \\
pi0.5 & 0.30 & 0.29 & 0.96 \\
Flow matching & 0.23 & 0.24 & 1.05 \\
DreamZero & 0.26 & 0.30 & 1.16 \\
FastWAM & 0.32 & 0.35 & 1.07 \\
\bottomrule
\end{tabular}

\end{table}

\textbf{To every model, the test demonstrations look like training ones.} A
test demonstration is between 0.90 and 1.16 times as far from its nearest
training demonstration as a training demonstration is from its own nearest
neighbour (Table~\ref{tab:recordings}). Each test demonstration has training
demonstrations that look about as similar as training demonstrations look to
each other, so the large gaps of Section~\ref{sec:compare} are not caused by
the test demonstrations looking unfamiliar.

\textbf{The session does contain one visual change, at demonstration 43.} ACT,
Diffusion and pi0.5 split the session into two blocks
(Figure~\ref{fig:recordings}). Splitting each encoder's features by camera
puts the boundary in the left arm's left fingertip (LLG) between demonstrations
42 and 43, for ACT, Diffusion, flow matching and pi0.5 alike; no other camera
changes at a common point. The untouched gel's image shows why: between those
two demonstrations the sensor's view shifted, moving the ridge and the pad's
edges and adding new marks on the gel (Appendix Figure~\ref{fig:sensorshift}).
The shift is a step, not a drift: every demonstration before it looks like the
earlier ones and every one after it like the later ones. It does not separate
test from training, since 15 training demonstrations (43--57) come after it
too. DreamZero and FastWAM barely show it, probably because our features for
them average each fingertip onto about four cells of the autoencoder's latent,
too coarse to see the ridge move.}

\rev{\subsection{Learning curves}
\label{sec:curves}

Figure~\ref{fig:curves} follows each model's training and test error over the
first ten steps through training (Section~\ref{sec:measures}). The last point
of each curve reproduces the model's row of Table~\ref{tab:actions} to within
what a rerun moves it: the curves come from fresh runs with the same settings,
and ACT's last point over the first ten steps is 8.7 against 7.9 in the table.

\begin{figure}[t]
  \centering
  \includegraphics[width=0.85\textwidth]{curves.png}
  \caption{Action error over the first ten steps on the training (dashed) and
  test (solid) demonstrations, at every saved checkpoint of a rerun with the
  same settings. The dotted line marks the checkpoint with the lowest test
  error.}
  \label{fig:curves}
\end{figure}

\textbf{The test error stops improving long before training ends; the training
error does not.} Diffusion's test error levels off at about 5.2 by
\num{40000} of its \num{100000} steps, while its training error keeps falling,
from 3.5 to 1.8. Flow matching's levels off by \num{20000} of \num{60000}.
ACT's training error falls from 8.4 to 2.5 over the whole run, while its test
error moves only from 10.5 to about 8, and over its whole plan not at all (14.7
at \num{5000} steps, 14.3 at the end). The large gaps of
Table~\ref{tab:actions} are therefore built late in training: the extra steps
fit the training demonstrations more closely without helping on new ones.
They do not make the test error worse either, so stopping earlier would save
time, not improve the models.

\textbf{pi0.5 never pulls apart.} Its training and test errors fall together
and level off at about \num{20000} steps, with a gap of 1.1 throughout: it
stops learning rather than overfitting, as Section~\ref{sec:compare} found.

\textbf{The trainer's own held-out loss would have misled us.} LeRobot also
reports each model's training loss on the test demonstrations. For Diffusion
it rises fourfold from \num{10000} steps (0.029 to 0.119), and for flow
matching it rises too, which reads as severe overfitting, yet neither model's
test action error gets worse. These losses score how well a model removes
noise from the recorded actions, not how close its plan is to them.
\pending{DreamZero and FastWAM learning curves.}}

\subsection{Cropping the tactile images}
\label{sec:crop}

Table~\ref{tab:crop} gives the test error of every model with and without each
crop.

\begin{table}[t]
  \centering
  \caption{Test action error (RMSE over each model's own plan) with each
  tactile crop (Section~\ref{sec:cropmethod}). \emph{Rows, unstretched} exists
  for ACT only. Each entry is one training run; \new{the lowest in each row is in
  bold.}}
  \label{tab:crop}
  \small
  \begin{tabular}{lrrrr}
\toprule
model & no crop & rows & rows, unstretched & four edges \\
\midrule
ACT & 14.09 & 14.68 & 14.08 & -- \\
Diffusion & 9.18 & 9.21 & n/a & -- \\
pi0.5 & 18.35 & 18.11 & n/a & -- \\
DreamZero & 7.02 & n/a & n/a & -- \\
FastWAM & 3.94 & n/a & n/a & -- \\
\bottomrule
\end{tabular}

\end{table}

\textbf{Cropping the tactile images changes almost nothing.} With the rows-only
crop, ACT, Diffusion and pi0.5 move by a few per cent at most, in both
directions. That is within what repeating a run with another random seed moves
Diffusion and pi0.5 (Appendix Table~\ref{tab:floor}). \new{The four-edge crop
leaves pi0.5 and Diffusion unchanged (18.34 against 18.35, 9.10 against 9.18)
and costs ACT as much as the rows-only crop does (14.70 against 14.68 and 14.09
uncropped).} \rev{FastWAM's four-edge crop leaves it unchanged too (3.99
against 3.94). The ridge crop, which also removes the gel's bright ridge, is no
different: pi0.5 18.49, Diffusion 9.06 and ACT 14.63, again within a few per
cent of uncropped and, for ACT, the same cost as its other crops. DreamZero's
ridge crop leaves it unchanged as well (6.98 against 7.02).}
\pending{DreamZero's four-edge crop and FastWAM's ridge crop.} For ACT, the one place the rows-only crop looked
harmful, over the far end of its plan, disappears when the image is not
stretched. So the stretching, not the loss of the rim, caused that small
penalty.

\textbf{The policies do not over-weight the edges of the tactile images.}
Appendix Figure~\ref{fig:rim} measures this directly with Grad-CAM. A value of 1 means the edge band
gets exactly its share of the weight by area; above 1, the model favours the
edge. On a test demonstration, every uncropped model sits between two thirds
and nine tenths of its share, at every band width. \new{So the strong edge
weighting of Figure~\ref{fig:gradcamevidence} happens at some moments but does
not hold on average.} Cropping moves the weight slightly towards the new edge,
partly for a mechanical reason: after cropping and stretching, the pixels
nearest the edge used to be further in.

\begin{figure}[t]
  \centering
  \includegraphics[width=0.6\textwidth]{patch_maps.png}
  \caption{\new{Patch occlusion maps: where in each fingertip image the plan
  comes from, averaged over 20 to 60 frames of test demonstration 58 (FastWAM:
  8 frames, $5\times5$ cells per sensor). Darker cells moved the plan more when
  painted over; each map is scaled to its own peak. A cropped model's map is
  drawn over the full image, so the white border is what it never sees. pi0.5
  reads only LLG and RLG.}}
  \label{fig:patchmaps}
\end{figure}

\new{\textbf{Cropping moves ACT's attention off the ridge; pi0.5's stays on it.}
Figure~\ref{fig:patchmaps} answers the question for models Grad-CAM cannot
reach, using patch occlusion. On each sensor we locate the five cells that move
the plan most. For uncropped ACT they sit 0.31 to 0.47 of the way across, around
the gel's ridge at about 0.3, and the outer fifth of the image gets 0.68 of its
share by area, in line with Grad-CAM. With the four-edge crop they move to 0.45
to 0.69, the right half of the gel, where the contact marks form. pi0.5 with the
four-edge crop still peaks at 0.33 to 0.35, on the ridge, which is what the
ridge crop removes. FastWAM's map is nearly flat: painting any cell moves its plan by at
least 40\,\% as much as the most important cell. Its pinned plan repeats
exactly, so this is the model and not sampling noise: it reacts to a change
anywhere in a fingertip image by about the same amount.}
\rev{Uncropped Diffusion peaks 0.37 to 0.53 across, just right of the ridge;
uncropped pi0.5 peaks in different places on its two sensors (0.61 and 0.23)
and its maps are flatter than the behaviour-cloning policies'; DreamZero's
peaks are scattered (0.13 to 0.73), and retrained flow matching's lie between
0.21 and 0.55. \textbf{No model over-weights the image edges:} the outer fifth
of the image gets 0.78 (Diffusion), 0.88 (DreamZero), 0.91 (pi0.5) and 0.95
(flow matching) of its share by area, all below one. Uncropped ACT, at 0.68, is
the most centred.}

\subsection{Which sensors the models use}
\label{sec:sensors}

We retrained ACT with three sets of inputs, all with proprioception: the
overhead camera alone; the overhead camera and one fingertip per arm (the pair
pi0.5 reads); and all five cameras. Table~\ref{tab:sensors} gives the result.
\new{The overhead-only model tests the reading of Appendix Figure~\ref{fig:framewise}: if
ACT uses touch only briefly, removing touch altogether should cost little.}

\begin{table}[t]
  \centering
  \caption{ACT trained on three sets of cameras, all with proprioception.
  Action error (RMSE) over its own 100-step plan and over the first ten steps.
  \new{The lowest error in each column is in bold.}}
  \label{tab:sensors}
  \small
  \begin{tabular}{lrrrrr}
\toprule
& \multicolumn{3}{c}{over its 100-step plan} & \multicolumn{2}{c}{first 10 steps} \\
\cmidrule(lr){2-4}\cmidrule(lr){5-6}
cameras & train & test & gap & test & gap \\
\midrule
overhead & -- & -- & -- & -- & -- \\
overhead + one fingertip per arm & -- & -- & -- & -- & -- \\
overhead + all four fingertips & 2.67 & 14.09 & 5.3 & 7.88 & 3.4 \\
\bottomrule
\end{tabular}

\end{table}

\new{\textbf{ACT generalises better without the fingertips.} Trained on the
overhead camera alone, it fits its training demonstrations exactly as well as
with all five cameras (2.62 against 2.67), but its test error is lower: 13.8
against 14.1 over its plan, and 5.8 against 7.9 over the first ten steps. So on
new demonstrations the fingertips do not help ACT, and may hurt it. Each setting
is one training run, and we have not measured how much a second run with
another seed would move these numbers (Appendix~\ref{sec:threats}).
Adding one fingertip per arm back does not help either: 14.8 over the plan and
7.5 over the first ten steps, between the other two. Neither tactile set beats
the overhead camera alone.}

\begin{figure}[t]
  \centering
  \includegraphics[width=0.5\textwidth]{stream_shares.png}
  \caption{How much each input moves each model's plan (input contribution,
  Section~\ref{sec:measures}). Each bar adds up to 100\,\%, so compare the
  order of inputs within a bar, not heights between bars. pi0.5 reads two
  fingertips; the others read four. \rev{DreamZero alone is given its past
  arm commands as an input.}}
  \label{fig:streams}
\end{figure}

Figure~\ref{fig:streams} shows which inputs the policies actually use. ACT
relies most on the overhead camera, then proprioception, with the fingertips
last. Diffusion and pi0.5 rely mostly on proprioception: about three quarters
of what moves pi0.5's plan is where the arms are. \new{FastWAM leans on
proprioception most of all, 86\,\%, with 11\,\% for the overhead camera and
3\,\% for the fingertips.} \rev{Retrained on whole images, flow matching
uses its cameras: proprioception falls to 75\,\% (98.7\,\% in the first run),
with 9\,\% for the overhead camera and 16\,\% for the fingertips, the largest
fingertip share of any model. DreamZero leans on proprioception most of all
(88\,\%), then on its past commands (10\,\%); all five cameras together move
its plan by only 2.3\,\%. The model that plans the next ten steps second best
barely reads its cameras to do it.}


\subsection{How well the world action models predict the cameras}
\label{sec:wm}

Table~\ref{tab:wm} scores both world action models on every sensor. The first
column is reconstruction, the frames the model saw passed through its own
autoencoder, which is the sharpest any of its predictions can be. \new{For
DreamZero this is its frozen image autoencoder, applied to each frame; for
FastWAM it is Wan2.2's video autoencoder, which compresses groups of four
frames together.} The next columns compare its predicted future frames with the
real ones and with repeating the last frame it saw.

\begin{table}[t]
  \centering
  \caption{World action models on the seven test demonstrations, per sensor, in
  PSNR (dB, higher is better). \emph{Reconstruction}: seen frames through the
  model's own autoencoder. \emph{Prediction}: predicted future frames.
  \new{\emph{Repeat last}: the last seen frame used as the prediction for every
  future time.} Both averaged over the prediction horizon. \emph{Margin}:
  prediction minus repeat last, at \SI{0.8}{\second} ahead and at the end of
  each model's horizon (\SI{4.8}{\second} for DreamZero, \SI{1.1}{\second} for
  FastWAM). Each model works at its own resolution, so compare margins, not
  PSNR, between models.}
  \label{tab:wm}
  \small
  \begin{tabular}{llrrrrr}
\toprule
& & & & & \multicolumn{2}{c}{margin} \\
\cmidrule(lr){6-7}
sensor & model & reconstruction & prediction & repeat last & at \SI{0.8}{\second} & at end \\
\midrule
overhead & DreamZero & 17.3 & 15.5 & 14.1 & -4.1 & +3.2 \\
LLG & DreamZero & 37.0 & 30.1 & 25.9 & -1.3 & +5.6 \\
LRG & DreamZero & 34.9 & 28.0 & 22.9 & -1.4 & +6.7 \\
RLG & DreamZero & 33.9 & 24.8 & 20.7 & +0.0 & +5.0 \\
RRG & DreamZero & 37.0 & 26.8 & 23.6 & +1.5 & +2.6 \\
overhead & FastWAM & 33.8 & 18.5 & 18.6 & +0.2 & +0.4 \\
LLG & FastWAM & 46.7 & 33.0 & 32.9 & +5.2 & +5.2 \\
LRG & FastWAM & 46.1 & 31.5 & 33.4 & +2.7 & +4.3 \\
RLG & FastWAM & 45.9 & 28.5 & 26.5 & +5.0 & +6.3 \\
RRG & FastWAM & 47.2 & 30.0 & 29.0 & +3.0 & +3.1 \\
\bottomrule
\end{tabular}

\end{table}

\textbf{Most of the prediction error comes from predicting, not from
compressing the images.} DreamZero's small image autoencoder is the exception
on the overhead camera: it reproduces a seen frame at only 17.3\,dB, and the
model's predictions reach 15.5\,dB, close to that ceiling. \new{(At 17.3\,dB a typical
pixel is off by 14\,\% of the full brightness range, a visibly blurred image;
at 35\,dB it would be off by under 2\,\%.)} On the fingertips its
autoencoder reaches 34 to 37\,dB, but its predictions only 25 to 30\,dB.
FastWAM's pretrained autoencoder is much sharper, 34\,dB on the overhead camera
and about 46\,dB on the fingertips, and its predictions sit 13 to 17\,dB below
that. So a better autoencoder would remove little of either model's error,
except DreamZero's on the overhead camera.

\textbf{DreamZero predicts the cameras better than repeating the last frame,
on every sensor, from \SI{1.6}{\second} ahead onwards.} Its margin is roughly 3
to 8\,dB, and largest on the fingertips, where repeating the last frame is
already strong because the gels change only at contact. At \SI{0.8}{\second} it
only ties, because the future still looks almost exactly like the present.
\new{In Figure~\ref{fig:filmstrip} its reconstruction of the seen frame is
already blurred, but it predicts the arms leaving and the folded shorts staying
where they are, as really happened.}

\textbf{FastWAM predicts sharp frames, but a frame-averaged score barely
rewards them.} It predicts only about a second ahead. On the overhead camera it
ties with repeating the last frame, although its frames move the grippers in
the right direction (Figure~\ref{fig:filmstrip}). \new{Its difference row
shows the error concentrated on the moving arms and cloth, where its motion lags
the real one; everywhere else it is flat grey.} The garment fills most of the
frame and hardly moves within a second, so repeating the last frame is already
right about most pixels. \textbf{A score averaged over the
whole frame can rate a correct prediction of a small moving part as no better
than doing nothing.} On the fingertips, where contact changes many pixels,
FastWAM starts below repeating the last frame and ends 3 to 6\,dB above it on
every sensor. At the one moment both models predict, \SI{0.8}{\second} ahead,
FastWAM is ahead on both panels of Appendix Figure~\ref{fig:predmargin}.

\begin{figure}[t]
  \centering
  \includegraphics[width=0.6\textwidth]{filmstrip_h2h.png}
  \caption{\new{FastWAM and DreamZero predicting from the same seen frame
  (test demonstration 58, 8\,s in, as the arms release the folded shorts), on
  one time axis. For the overhead camera and one fingertip (LLG) the rows are
  what happened, each model's prediction, and its difference from what happened
  (truth minus prediction, rescaled so that mid-grey is no error). The
  \emph{seen} column shows each model's reconstruction of the frame it was
  given. Blank cells are times a model does not predict. Each model is shown at
  its own resolution. Appendix Figure~\ref{fig:filmstripearly} shows an earlier
  moment.}}
  \label{fig:filmstrip}
\end{figure}

\rev{\subsection{Does FastWAM's touch prediction follow the gripper?}
\label{sec:grasp}

We scored FastWAM's fingertip predictions over the 36 grasps of the test
demonstrations (35 distinct windows, since two grasps start on the same frame),
on the two fingertips of the grasping arm. Table~\ref{tab:grasp} gives the mean
absolute pixel error as a percentage of the pixel range, and
Figure~\ref{fig:grasp} shows one grasp.

\begin{table}[t]
  \centering
  \caption{FastWAM around the 36 test grasps: mean absolute pixel error (\% of
  the pixel range) of the grasping arm's two fingertips, over the eight
  predicted frames. \emph{Change from recorded} is how far the prediction moves
  when the current gripper position given to the model is changed from the
  recorded (open) value to closed. \emph{Repeat-last error} is the error of
  repeating the seen frame, which is also how much the gel really changes.}
  \label{tab:grasp}
  \small
  \begin{tabular}{lrrr}
\toprule
condition & error & change from recorded & repeat-last error \\
\midrule
recorded & 5.44 & 0.00 & 8.00 \\
closed & 5.47 & 1.38 & 8.00 \\
\bottomrule
\end{tabular}

\end{table}

\begin{figure}[t]
  \centering
  \includegraphics[width=\textwidth]{grasp_filmstrip_state.png}
  \caption{One grasp (test demonstration 60, left arm, fingertip LLG). Rows:
  what happened; FastWAM's prediction with the recorded gripper position;
  its prediction when told the gripper is already closed; and the difference
  between the two predictions (mid-grey is none). The \emph{seen} column shows
  the model's reconstruction of the frame it was given. The gripper is
  commanded closed at \SI{0.5}{\second}; the gel darkens at contact about a
  tenth of a second later.}
  \label{fig:grasp}
\end{figure}

\textbf{FastWAM predicts the contact, and when it happens, from the image it
sees.} In Figure~\ref{fig:grasp} it darkens the gel at the same moment the real
gel darkens, about a tenth of a second after the closing command, and it never
saw that command. Over all 36 grasps its error (5.4\,\%) is a third lower than
repeating the seen frame (8.0\,\%).

\textbf{Telling it the gripper is already closed changes almost nothing.} The
prediction moves by 1.4\,\% when the gripper position is set to closed, against
8.0\,\% of real change in the gel, and its error is the same (5.5\,\%). The
difference row of Figure~\ref{fig:grasp} is flat grey. FastWAM reads the
coming grasp from the images (where the arm is, where the cloth is), not from
the gripper's position. \pending{the FastWAM whose video reads the actions,
with the grasp commanded and with the gripper held open.}}
